\begin{notebox}{Indexes and Tables of Content}
A brief Table of Contents is provided at the beginning of this document (see page \pageref{label:short_toc}), outlining only the chapters and main sections. 

A comprehensive Table of Contents can be found at the end of the thesis (page \pageref{label:full_toc}), alongside the List of Figures (page \pageref{label:list_of_figures}) and the Bibliography (page \pageref{label:bib_references}). 
\end{notebox}

\begin{notebox}{Source code and repository structure}
Chapters \ref{chap:implementation_and_simulation} and \ref{chap:simulating_2d_color_codes} show the result of multiple simulations and code execution. 

Small code snippets are shown to illustrate how the different tools and libraries work. The complete code required to execute the simulations presented in this work is available in a git repository at \url{https://github.com/othermore/TFM_QEC_2D_color_code_decoders}. \textbf{The line numbers in the original files for each snippet are displayed} to facilitate locating them in the source files.

The code is located in the \texttt{src} directory. Each chapter has its own directory and python notebook. The specific code for each simulation is located in separate files (starting with \texttt{s} and the number of the simulation) and included in the notebook. 

Libraries and resources required for multiple simulations within a given chapter, are included in the chapter's directory.
\end{notebox}